% ৩.১৪ মৌলিক গেট
\section{মৌলিক গেট: AND, OR, NOT}

\subsection{AND Gate}
\begin{itemize}
\item Symbol, Boolean: $Y=A\cdot B$, Truth table
\end{itemize}

\subsection{OR Gate}
\begin{itemize}
\item Symbol, Boolean: $Y=A+B$, Truth table
\end{itemize}

\subsection{NOT Gate}
\begin{itemize}
\item Symbol, Boolean: $Y=\overline{A}$, Truth table
\end{itemize}

\begin{learningbox}
এই টপিক শেষে শিক্ষার্থী পারবে:
\begin{itemize}
\item AND, OR ও NOT gate-এর symbol, Boolean expression ও truth table লিখতে।
\item মৌলিক gate দিয়ে simple logic circuit ব্যাখ্যা করতে।
\item বাস্তব switch analogy ব্যবহার করে gate-এর output বুঝতে।
\end{itemize}
\end{learningbox}

\begin{center}
\fbox{\parbox{0.9\textwidth}{
\centering
\textbf{Logic gate symbol placeholders}\\[4pt]
AND gate: দুই input বাম পাশে, output ডান পাশে।\\
OR gate: curved input side, output ডান পাশে।\\
NOT gate: triangle with small bubble at output।
}}
\end{center}

\begin{center}
\begin{tabular}{|c|c|c|c|}
\hline
$A$ & $B$ & $A\cdot B$ & $A+B$ \\
\hline
0 & 0 & 0 & 0 \\
\hline
0 & 1 & 0 & 1 \\
\hline
1 & 0 & 0 & 1 \\
\hline
1 & 1 & 1 & 1 \\
\hline
\end{tabular}
\end{center}

\begin{center}
\begin{tabular}{|c|c|}
\hline
$A$ & $\overline{A}$ \\
\hline
0 & 1 \\
\hline
1 & 0 \\
\hline
\end{tabular}
\end{center}

\begin{examplebox}{Switch analogy}
AND gate হলো series switch-এর মতো: দুই switch on না হলে current যাবে না। OR gate হলো parallel switch-এর মতো: যেকোনো একটি switch on হলেই current যেতে পারে।
\end{examplebox}
